\subsection{Isotopy extension theorems}

Smooth isotopies are special kinds of homotopies, and as such we can define similar concepts.

\bdefn

    In general an isotopy is a smooth map $\psi_\bul\colon[0,1]\times X\to Y$ such that $\psi_t\colon X\to Y$ is a diffeomorphism for all $t$.
    We may also require only that $\psi_t$ be an embedding for all $t$; the meaning will be made clear by context.
    If $A\subseteq X$ is an arbitrary subset, an {\emphcolor isotopy relative to $A$} is an isotopy $\psi_\bul$ such that $\psi_t|_A=\psi_0|_A$ for all $t$.
    If $Q\subseteq X$ is a submanifold, an {\emphcolor isotopy of $Q$} is a smooth isotopy (of embeddings) $\psi_t$ such that $\psi_0$ is the inclusion $Q\inj X$.

    All of these notions can be specialized to {\emphcolor symplectic isotopies}, where $\psi_t\colon Q\to M$ is a symplectomorphism.

\edefn

In the following for real $r>0$, $B(r)$ denotes the closed ball of radius $r$ centered at $0$ in $\bR^{2n}$.
This inherits the standard symplectic structure from $\bR^{2n}$.

\bthrm[name={sympballext}]

    Let $(M,\omega)$ be a connected symplectic manifold without boundary.
    \benum
        \item Let $\lambda>0$, then every symplectic embedding $\psi\colon B(r)\to M$ extends to a symplectic embedding $B(r)\inj M$ for some $r>\lambda$.
        \item Let $0\leq\lambda<r$ and $\psi_0,\psi_1\colon B(r)\to M$ be symplectic embeddings which are connected by a smooth family of
        symplectic embeddngs $\psi_t\colon B(\lambda)\to M$.
        Then there is a Hamiltonian isotopy of $M$ $\phi_t$ and a constant $\lambda<\rho<r$ such that
        $$ \phi_0 = \id,\qquad \phi_1\circ\psi_0|_{B(\rho)} = \psi_1|_{B(\rho)},\qquad \phi_t\circ\psi_0|_{B(\lambda)} = \psi_t . $$
        Thus $\phi_t\circ\psi_0$ is an extension of the isotopy $\psi_t$ onto the larger ball $B(\rho)$.
    \eenum

\ethrm

We will first need the following lemma.

\blemm

    Let $\psi_\bul\colon[0,1]\times B(\lambda)\to M$ be a symplectic isotopy (of embeddings).
    Then there is a Hamiltonian isotopy $\phi_\bul\colon[0,1]\times M\to M$ such that
    $$ \phi_0 = \id,\qquad \phi_t\circ\psi_0 = \psi_t . $$

\elemm

\bproof

    Let $X_t$ be the time-dependent vector field generated by $\psi_t$.
    We have that $\psi_t^*\omega=\omega_0$ and thus we differentiate
    $$ 0 = \drv t\psi_t^*\omega = \psi_t^*\L_{X_t}\omega = \psi_t^*\d\iota(X_t)\omega = \psi_t^*\d\omega(X_t,\bul)
    = \d\omega\parens{\partial_t\psi_t,\d\psi_t\bul} . $$
    Thus we have a smooth family of closed $1$-forms
    $$ \alpha_t = \omega(\partial_t\psi_t,\d\psi_t\bul) \in \Omega^1(B(\lambda)) . $$
    Because $B(\lambda)$ is contractible, this means $\alpha_t$ is exact for all $t$.
    So there is a smooth family of functions $h_t\colon B(\lambda)\to\bR$ such that
    $$ \d h_t = \alpha_t,\qquad h_t(0) = 0 . $$

    Note that we can extend $\psi_t$ to an isotopy of embeddings $\psi_\bul\colon[0,1]\times\bR^{2n}\to M$.
    Then via this lift $h_t$ to $H_t\colon M\to\bR$ such that
    $$ H_t\circ\psi_t = h_t . $$
    Then consider the Hamiltonian isotopy $\phi_t$ generated by $H_t$:
    $$ \drv t\phi_t = X_{H_t}\circ\phi_t,\qquad \phi_0 = \id,\qquad \iota(X_{H_t})\omega = \d H_t . $$
    Now notice that
    $$ \omega(\partial_t(\phi_t\circ\psi_0),\d\psi_t\bul) = \omega(X_{H_t}\circ\phi_t\circ\psi_0,\d\psi_t\bul) = \d H_t(\d\psi_t\bul) = \d(H_t\circ\psi_t) = \d h_t
    = \omega(\partial_t\psi_t,\d\psi_t\bul) . $$
    Because $\psi_t$ is symplectic and $\omega$ nondegenerate, this means $\partial_t(\phi_t\circ\psi_0)=\partial_t\psi_t$.
    And because they agree at $t=0$ we have $\phi_t\circ\psi_0=\psi_t$.
    \qed

\eproof

\bproof[\refmath{sympballext} (1)]

    The symplectic embedding $\psi\colon B(\lambda)\to M$ extends to a smooth embedding $\psi\colon B(\lambda+\epsilon)\to M$.
    Then the $2$-form
    $$ \psi^*\omega \in \Omega^2(B(\lambda+\epsilon)) $$
    agrees with $\omega_0$ on $B(\lambda)$.
    Because $B(\lambda),B(\lambda+\epsilon)$ are contractible, by homotopy invariance of de Rham cohomology, because $\omega_0-\psi^*\omega$ is zero on $B(\lambda)$
    and thus exact, $\omega_0-\psi^*\omega$ is exact on $B(\lambda+\epsilon)$.
    So let $\alpha\in\Omega^1(B(\lambda+\epsilon))$ such that
    $$ \alpha|_{B(\lambda)} = 0,\qquad \d\alpha = \psi^*\omega - \omega_0 . $$
    By potentially shrinking $\epsilon$, we may assume $\omega_s=\omega_0+s\d\alpha$ is nondegenerate for all $0\leq s\leq1$.
    By \refmath[lemma]{mosercompactisotopy} there is an isotopy $\chi_t\in\Diff(B(\lambda+\epsilon))$ such that
    $$ \chi_0 = \id,\qquad \chi_s|_{B(\lambda)} = \id,\qquad \chi_s^*\omega_s = \omega_0 . $$
    Thus $\chi_1^*\psi^*\omega=\omega_0$ and thus $\psi\circ\chi_1$ is the desired extension.
    \qed

\eproof

\bproof[\refmath{sympballext} (2)]

    We first consider $\lambda>0$.
    \TODO{This?}

\eproof

\TODO{Continue this; just 3.3.2}

